% !PRINTED 117
\weberpage{117}

Wir geben der $\sqrt R$ eines der beiden Zeichen und erhalten nach (5)
\begin{equation}
\begin{array}{rclcrcl}
u^3&=&-\dfrac b2+\sqrt R,&& v^3&=&-\dfrac b2-\sqrt R,\\[0.75em]
u&=&\sqrt[3]{-\dfrac b2+\sqrt R},&& v&=&\sqrt[3]{-\dfrac b2-\sqrt R},
\end{array}
\tag{8}
\end{equation}
und also
\begin{equation}
y=\sqrt[3]{-\frac b2+\sqrt R}+\sqrt[3]{-\frac b2-\sqrt R}.
\tag{9}
\end{equation}

Die Multiplication der beiden Ausdrücke (8) ergiebt
\[
 u^3v^3=-\frac{a^3}{27}
\]
und zeigt, dass, wenn für $u$ ein bestimmter unter den drei Werthen der Cubikwurzel genommen wird,
\begin{equation}
v=-\frac{a}{3u}=-\frac{a}{3\sqrt[3]{-\dfrac b2+\sqrt R}}
\tag{10}
\end{equation}
jedenfalls einer der drei Werthe von $\sqrt[3]{-\dfrac b2-\sqrt R}$ ist. Aus (4) folgt aber, dass, nur wenn wir diesen Werth für $v$ nehmen, die Gleichung (1) durch $y=u+v$ befriedigt ist. Wir erhalten so, den drei Werthen von $u$ entsprechend, drei Wurzeln der Gleichung (1). Aus der Existenz einer Wurzel der cubischen Gleichung ergiebt sich aber auch daraus schon die Existenz von dreien, dass bereits nachgewiesen ist, dass die quadratische Gleichung zwei Wurzeln hat.

Um aus der einen Wurzel (9) die beiden anderen abzuleiten, setzen wir
\[
 \alpha=u+v
\]
und dividiren $y^3+ay+b$ durch $y-\alpha$. Der Quotient, dessen Wurzeln die beiden anderen Wurzeln $\beta,\gamma$ der cubischen Gleichung (1) sind, ist
\[
 y^2+\alpha y+\alpha^2+a=0,
\]
woraus man nach (4) des vorigen Paragraphen
\[
 y=\frac{-\alpha\pm\sqrt{-3\alpha^2-4a}}{2}
\]
findet. Setzt man hierin $\alpha=u+v$, und nach (4) $a=-3uv$, so folgt
%CONT>>
